% Chapter 1: Introduction — Handbook 8.2.2
% ~1000 words

\section{Introduction}
\label{sec:introduction}

\subsection{Problem Statement}
\label{sec:intro-problem}

Since the inception of the digital age, many established industries have moved
online.  Auction systems in particular have seen a massive surge in popularity
due to the low barrier to entry for both buyers and sellers: eBay sold
\$7.2~million worth of goods within a year of launch~\cite{ebay_history2025}
and today hosts 2.5~billion live listings across 135~million active
buyers worldwide~\cite{ebay_about2025}, and this is only one platform.  However, their
security models have not kept pace.  The move from physical sealed envelopes
to digital platforms introduced new attack surfaces, including bid visibility,
platform trust, and metadata exposure. Existing systems address
these concerns incompletely, if at all.

Popular online marketplaces follow three distinct patterns.  eBay uses open
proxy bidding: the current highest bid is visible throughout the auction and
the platform operator sees all bids~\cite{ebay_help_bidding}.  Facebook
Marketplace is a listing-plus-messaging workflow with no formal auction
protocol and weak buyer protections~\cite{fb_marketplace_help}.  Blockchain
marketplaces such as OpenSea~\cite{opensea2025} use English style auctions where
every bid is an on-chain transaction, publicly visible on all supported
chains~\cite{opensea_chains2025} and permanently recorded solving the
trusted-auctioneer problem via smart contracts, but providing zero bid privacy.

Government procurement offers a contrasting model.  Under the US Federal
Acquisition Regulation, sealed bidding keeps all bids confidential until a
public opening event~\cite{far_part14}.  Bids are private \emph{during}
submission, but a trusted auctioneer (the contracting officer) sees every bid
at opening and determines the winner.  Sealed bidding thus protects bids from
\emph{other bidders} but not from the \emph{auctioneer}.

Daian et al.\ demonstrated that even blockchain-based commit-reveal schemes
are vulnerable to front-running and miner-extractable value
(MEV)~\cite{daian2019_flashboys2}, where transaction ordering can be
exploited during the reveal phase.  Cryptographic sealed-bid auctions can
hide non-winning bids entirely~\cite{sako2000_hide_losers}, and multi-party
computation (MPC) can remove the need for a trusted auctioneer
altogether~\cite{bogetoft2009_mpc_live}.  However, deploying MPC over a
transport layer that also hides participants' network identities remains
largely unexplored.

\begin{table}[H]
\centering
\small
\renewcommand{\arraystretch}{1.3}
\makebox[\textwidth][c]{%
\begin{tabular}{|l|l|l|l|l|}
\hline
\textbf{System} & \textbf{Bids visible?} & \textbf{Trusted auctioneer?} & \textbf{Decentralised?} & \textbf{Bidder identity} \\
\hline
eBay              & Yes (open ascending) & Yes           & No  & Real name + address \\
\hline
OpenSea           & Yes (on-chain)       & No (contract) & Yes & Pseudonymous (wallet, IP visible) \\
\hline
US govt.\ sealed bid & No (until opening) & Yes (sees all)& No  & Real name (registered) \\
\hline
\textbf{Lattice(This project)} & \textbf{No (MPC)} & \textbf{No} & \textbf{Yes} & \textbf{Pseudonymous (key, IP hidden)} \\
\hline
\end{tabular}}%
\caption{Comparison of auction privacy models~\cite{ebay_help_bidding,opensea2025,far_part14}.}
\label{tab:auction-comparison}
\end{table}

\subsection{Security Policy}
\label{sec:intro-security-policy}

Privacy in this system meant protecting specific information from specific
adversaries.  Table~\ref{tab:security-policy-intro} summarises the security
policy; Chapter~\ref{sec:context} expands on the threat model in the context
of Veilid's network-layer protections.

\begin{table}[H]
\centering
\small
\renewcommand{\arraystretch}{1.3}
\makebox[\textwidth][c]{%
\begin{tabular}{|l|l|l|}
\hline
\textbf{What is protected} & \textbf{From whom} & \textbf{Mechanism} \\
\hline
Non-winning bid values & All parties (incl.\ seller) & MPC: only the maximum is revealed to party~0 \\
\hline
Bidder network identity (IP) & Network observers, other bidders & Veilid private routes (onion-routed, multi-hop) \\
\hline
Listing content before purchase & Non-winners & AES-256-GCM; key released only to verified winner \\
\hline
Bid existence and timing & Partially protected & DHT participation metadata remains visible \\
\hline
\end{tabular}}%
\caption{Security policy: what is protected and from whom.}
\label{tab:security-policy-intro}
\end{table}

The \emph{adversary model} considered four types, summarised in
Table~\ref{tab:adversary-model}.

\begin{table}[H]
\centering
\small
\renewcommand{\arraystretch}{1.3}
\makebox[\textwidth][c]{%
\begin{tabular}{|l|l|l|}
\hline
\textbf{Adversary} & \textbf{Capability} & \textbf{Goal} \\
\hline
Honest-but-curious bidder & Follows protocol, observes all available data & Learn other bids or bidder identities \\
\hline
Malicious seller & Controls listing, party~0 output, key release & Grief (withhold key), collude with preferred bidder \\
\hline
Passive network observer & Monitors DHT writes and message patterns & Deanonymise participants \\
\hline
Colluding parties & Multiple participants share private state & Break MPC threshold, reconstruct secret inputs \\
\hline
Dishonest participant & Participates in auction, observes protocol & Deanonymise other bidders or link pseudonyms to IPs \\
\hline
\end{tabular}}%
\caption{Adversary model.}
\label{tab:adversary-model}
\end{table}

Key \emph{trust and design assumptions}:

\begin{itemize}[noitemsep]
  \item The seller is semi-trusted: as MPC party~0 they learn the winning bid
    value and receive an ephemeral onion route to the winner, but not the
    winner's network identity.  The seller controls decryption key release
    and can grief by withholding the key.
  \item Veilid's private routing provides transport-layer anonymity; the
    effective anonymity set depends on both the number of relay nodes
    available in the overlay and the number of hops configured per
    private route. In this project's devnet; due to the devnet size, the hops default to 1.  More relay nodes increase the pool
    of possible intermediaries; more hops make traffic analysis harder
    but add latency per round trip.
  \item Liveness requires all MPC parties to remain online throughout
    computation.  A single party going offline aborts the auction.
  \item The MPC security guarantee depends on the protocol backend:
    MASCOT~\cite{keller2016_mascot} tolerates up to $N{-}1$ corrupt parties
    (dishonest majority)- and this is justified by the seller getting extra information in comparison to the bidding participants.
    Shamir~\cite{shamir1979_secret_sharing} requires
    an honest majority of at least $\lfloor N/2 \rfloor + 1$ honest parties.
  \item The system operates on a private devnet, not the public Veilid
    network.  All nodes are controlled by the tester; no real-world
    adversaries or network conditions are present or have been tested.
  \item Each Veilid node is assumed to have a unique public IP address.
    The devnet simulates this via LD\_PRELOAD IP spoofing
    (Section~\ref{sec:ctx-veilid-bootstrap}).
  \item Veilid supports only one active private route per node at a time.
    MPC sessions are therefore sequential; concurrent auctions on the same
    node are not possible.
  \item Party assignment is deterministic (sorted by bid timestamp).
    Any deterministic assignment from public data (timestamps, public keys,
    or otherwise) is publicly derivable and constitutes a known side channel.
    With MASCOT this is low severity, as the protocol is secure regardless
    of party identity; with Shamir it is more concerning if an attacker
    can target specific parties for corruption
    (Section~\ref{sec:results-security}).
  \item No Sybil resistance: Veilid itself provides no defence against
    a single adversary creating multiple pseudonymous identities.  By
    extension, nothing prevents one party from placing multiple bids in
    the same auction under different keys.  Mitigating this would require
    an identity layer that is out of scope.
  \item The auction is first-price only: the highest bidder wins and pays
    their bid. 
  \item No payment or settlement layer is integrated.  The auction ends
    at decryption key delivery; actual exchange of goods or funds is
    out of scope.
\end{itemize}

\subsection{Project Objectives}
\label{sec:intro-objectives}

The main objective was to deliver a working peer-to-peer marketplace on the
Veilid network supporting sealed-bid listings with private content unlock for
the winning bidder.  Four sub-objectives were defined in the PDD
(Appendix~\ref{app:pdd}):

\begin{enumerate}
  \item \textbf{Listings and purchases.}  End-to-end auction flow: create
    listing $\rightarrow$ submit bid $\rightarrow$ determine winner
    $\rightarrow$ complete purchase.
  \item \textbf{Sealed-bid via MPC.}  Integrate a multi-party computation
    protocol to select the highest valid bid without revealing non-winning bids.
  \item \textbf{MPC-gated decryption.}  Bind winner selection to content
    access: only the verified winner obtains the decryption key.
  \item \textbf{Settlement on Monero testnet.}  Simulate wallet
    initialisation, deposits, and release on a private Monero testnet.
    This objective was descoped as it was not inline with the research
    goals of the project, being more suited to a production system
    (Section~\ref{sec:conc-objectives}).
\end{enumerate}

\subsection{Beneficiaries}
\label{sec:intro-beneficiaries}

The project consultant gained a proof-of-concept demonstrating Veilid's
capabilities for a non-trivial use case.  The Veilid development community
benefited from a real application exercising the framework's private routing
and DHT APIs.  The project also contributed empirical data on the overhead of
tunnelling MPC through a privacy-preserving overlay. Also contributed was the 
playground infrastructure to the Veilid community; as a seperate testing environment
for users to test their own applications in a ``sandbox''.

\subsection{Report Structure}
\label{sec:intro-structure}

Chapter~\ref{sec:output-summary} summarises the deliverables.
Chapter~\ref{sec:context} provides technical background on the building
blocks: Veilid, MP-SPDZ, and the stigmerge coordination pattern, and expands
the security policy into a full threat model.
Chapter~\ref{sec:method} describes the architecture, protocol design, and
testing strategy.  Chapter~\ref{sec:results} presents verification:
functionality, performance benchmarks comparing MASCOT and Shamir over Veilid,
and a security analysis of residual information leakage.
Chapter~\ref{sec:conclusions} gives a critical overview, reflections, and
future work.


